\documentclass[11pt]{article}
\usepackage{amsmath,amssymb,amsthm,mathtools,fullpage}
\newtheorem{theorem}{Theorem}[section]
\newtheorem{proposition}[theorem]{Proposition}
\newtheorem{lemma}[theorem]{Lemma}
\newtheorem{corollary}[theorem]{Corollary}
\theoremstyle{definition}
\newtheorem{definition}[theorem]{Definition}
\newtheorem{hypothesis}[theorem]{Hypothesis}
\theoremstyle{remark}
\newtheorem{remark}[theorem]{Remark}
\newcommand{\Z}{\mathbb{Z}}
\newcommand{\bl}{\boldsymbol{\lambda}}
\newcommand{\bs}{\mathbf{s}}
\newcommand{\bk}{\mathbf{k}}
\newcommand{\ket}[1]{\lvert #1 \rangle}
\newcommand{\bra}[1]{\langle #1 \rvert}
\title{Q63, gap 3: the wedge/Fock normalisation is trivial,\\
and what the $q$-power actually was}
\author{Clio}
\date{1 September 2026}

\begin{document}
\maketitle

\begin{abstract}
The Q63 paper of 31 August left as its blocking gap the claim that the identification
$u_{\bk}\leftrightarrow\ket{\bl,\bs}$ between Uglov's wedge basis and the combinatorial
Fock basis ``carries an unidentified $q$-power'': $[e_i,B_{-1}]$ was nonzero on $13$ of
$126$ test cases. The session brief proposed hypothesis (H4), that the power is the
cross-runner tie term $\tau$, to be tested against one named rival.

Both are wrong, and the reason is better than either. Uglov states outright
(\texttt{math/9905196}, the sentence after (4.6)--(4.7)) that the wedge action of
$U'_q(\widehat{\mathfrak{sl}}_n)$ in the multipartition indexation is \emph{identical}
to the combinatorial one, and that $U'_q(\widehat{\mathfrak{sl}}_n)$ and the Heisenberg
$\mathcal H$ commute. So there is no normalisation power, and $[e_i,B_{-1}]=0$ on the
nose. The observed failures were a defect in our level-$\ell$ straightening: it discarded
any wedge with a repeated index \emph{anywhere}, whereas (R1) licenses that only for an
\emph{adjacent} repeat. Level $1$ never exercises the difference, which is why the
implementation passed a $201$-block, $818$-coefficient answer-key match and still failed
at $\ell\ge2$. After the repair, $[e_i,B_{-1}]=0$ on $523/523$ and $[f_i,B_{-1}]=0$ on
$523/523$; the three mirror conventions fail on $76$--$92$ of $96$, so the check
discriminates.

Gap 3 then closes with more than a null result. We prove a closed form
(Theorem~\ref{thm:closed}) for every entry of $B_{-1}$ between multipartitions differing
in a single component:
$\bra{\nu}B_{-1}\ket{\bl}=(-q^{-1})^{h}q^{a}$, with $h$ the ribbon height and
\[
a\;=\;\#\{d'>d:\ \kappa\in M_{d'}\}\;+\;\#\{d'<d:\ \kappa+e\in M_{d'}\},
\]
a cross-runner count at the two positions the moving bead occupies. So the $q$-power that
gap~3 recorded is real, is bounded by $\ell-1$ for a proved reason, and is \emph{not
constant on a runner} for a proved reason --- but it belongs to the operator, not to a
change of basis. (H4) had the right species of object and the wrong home for it.

The engine is a locality lemma (Lemma~\ref{lem:local}): straightening the term
$k_j\mapsto k_j+N$ never disturbs a bead outside the closed window $[k_j,k_j+N]$, whose
length is exactly $N$. Two corollaries of the 31 August paper are thereby promoted from
\texttt{computed} to \texttt{proved}: every entry of $B_{-1}$ changing two or more
components is divisible by $(q-q^{-1})$, and $B_{-1}\vert_{q=1}$ is exactly the naive
per-runner ribbon operator.
\end{abstract}

\section{What was asked, and what the answer turned out to be}

The brief set (H4) --- \emph{the unidentified $q$-power is the tie term} --- and required
it be tested against the rival $a=\#\{\text{occupied runners}\}-1$, on the ground that a
check whose passing set is every object of that shape is not a check. That discipline was
right and it is what produced this session's result, but not in the way it anticipated:
the first thing to do with a quantity that is supposed to be unidentified is to look for
it in the source.

Uglov constructs the level-$\ell$ Fock space twice: combinatorially, as
$\mathcal F_{\bs}=\bigoplus_{\bl}\mathbb K\ket{\bl,\bs}$ with the
Jimbo--Misra--Miwa--Okado action (his Theorem 2.1, quoted from \cite{J,Fo}), and inside
the semi-infinite wedge $\Lambda^{s}$, where $\ket{\bl,\bs}:=u_{\bk}$ under the
indexation recalled in \S\ref{sec:conv}. Having written the wedge action in the
multipartition indexation he says, verbatim:

\begin{quote}
``These actions, and the actions of the generators $t_i$, $d$ \dots\ are identical with
those defined on the combinatorial Fock space $\mathcal F_{\bs}$.''
\end{quote}

\noindent and, two pages later,

\begin{quote}
``The actions of $U'_q(\widehat{\mathfrak{sl}}_n)$, $U'_p(\widehat{\mathfrak{sl}}_\ell)$
and $\mathcal H$ on $\Lambda^{s}$ are pairwise mutually commutative.''
\end{quote}

Together these say: the normalisation constant is $1$, and $[e_i,B_{-1}]=0$ identically.
Gap 3's premise --- that the identification carries an unidentified $q$-power --- is
therefore false, and the $13$ failures out of $126$ were a defect in our own pipeline.
\S\ref{sec:bug} locates it. What survives, and is worth the session, is
\S\ref{sec:closed}: the $q$-power that gap 3 \emph{measured} is a real feature of
$B_{-1}$, and it has an exact closed form.

\section{Conventions}\label{sec:conv}

Fix the rank $n=e\ge2$ and the level $\ell\ge1$; put $N=n\ell$.

Uglov's wedge basis is $u_{\bk}=u_{k_1}\wedge u_{k_2}\wedge\cdots$ for
$k_1>k_2>\cdots$, stabilising to $k_i=s-i+1$. Each $k\in\Z$ factorises uniquely as
\begin{equation}\label{eq:kcd}
k \;=\; c + n(d-1) - N\mu,\qquad c\in\{1,\dots,n\},\ d\in\{1,\dots,\ell\},\ \mu\in\Z,
\end{equation}
and the bead $k$ is read as sitting on \emph{runner} $d$ at position $\kappa=c-n\mu$.
Collecting the runner-$d$ positions in decreasing order as $\kappa^{(d)}_1>\kappa^{(d)}_2>\cdots$
and setting $\lambda^{(d)}_i=\kappa^{(d)}_i-s_d+i-1$ recovers the multipartition
$\bl$ and multicharge $\bs$; equivalently
\[
M_d(\bl,\bs)=\{\lambda^{(d)}_j-j+1+s_d\ :\ j\ge1\},
\]
the runner-$d$ Maya set. We write $M\subset\Z$ for the bead set in the $k$-coordinate.
The pair $(c,d)$ will be called the \emph{label} of $k$.

Two facts from \eqref{eq:kcd} are used repeatedly and neither needs a computation.
First, $\kappa$ and $\kappa+n$ have the same $c$ and the same $d$; so the Heisenberg
shift $k\mapsto k+N$ is $\kappa\mapsto\kappa+e$ on a single runner, and it preserves the
label. Second, the $N-1$ integers strictly between $k$ and $k+N$ realise every residue
modulo $N$ except that of $k$; hence each label other than $k$'s own occurs there
\emph{exactly once}.

\begin{definition}[Chevalley action, Uglov Thm.\ 2.1]\label{def:chev}
For a node $\gamma=(i,j,b)$ of $\bl$ put $x(\gamma)=s_b+j-i$ and
$\mathrm{res}(\gamma)=x(\gamma)\bmod n$. Order nodes by
$\gamma<\gamma'$ iff $x(\gamma)<x(\gamma')$, or $x(\gamma)=x(\gamma')$ and $b<b'$.
With $N^{<}_i(\bl,\boldsymbol\mu)$ the number of addable $i$-nodes of $\bl$ below
$\gamma$ minus the number of removable ones,
\[
e_i\ket{\boldsymbol\mu,\bs}=\sum q^{-N^{<}_i(\bl,\boldsymbol\mu)}\ket{\bl,\bs},
\qquad
f_i\ket{\bl,\bs}=\sum q^{+N^{>}_i(\bl,\boldsymbol\mu)}\ket{\boldsymbol\mu,\bs}.
\]
\end{definition}

\begin{remark}[The tie-break is not pinnable at $\ell=1$]\label{rem:tie}
The 31 August paper wrote that the tie-break direction and the choice of side ``were
fixed by requiring that the $\ell=1$ specialisation reproduce the proved Q59 theorem''.
That is true of the side and false of the tie-break: at $\ell=1$ no two nodes of distinct
components compete for a shifted content, so the tie-break is invisible. It is pinned
here, by the source, as $b<b'\Rightarrow\gamma<\gamma'$; and \S\ref{sec:comm} shows that
the mirror choice is detected --- and rejected --- by $[e_i,B_{-1}]=0$. It is also worth
recording what does \emph{not} pin it: the $\mathfrak{sl}_2$ relation
$[e_i,f_i]=[N_i(\bl)]_q$ holds for \emph{both} tie-breaks and both sides
($534/534$ checks each, three conventions). A relation satisfied by every convention is
not a convention check.
\end{remark}

\begin{definition}[Heisenberg, Uglov (3.13)]\label{def:B}
$B_m(u_{\bk})=\sum_{j} u_{k_1}\wedge\cdots\wedge u_{k_j-Nm}\wedge\cdots$; in particular
$B_{-1}$ replaces one $k_j$ by $k_j+N$ and then straightens.
\end{definition}

\begin{definition}[Straightening, Uglov Prop.\ 3.16]\label{def:R}
For $k_1\le k_2$ with labels $(c_1,d_1),(c_2,d_2)$ set $\gamma=(c_2-c_1)\bmod N$ and
$\delta=n(d_2-d_1)\bmod N$. Then $u_{k_1}\wedge u_{k_2}$ equals
\[
\varepsilon\,u_{k_2}\wedge u_{k_1}\;+\;\sum_{z>0} c_z\,u_{k_2-z}\wedge u_{k_1+z},
\qquad
\varepsilon=\begin{cases}
-1 & \gamma=0,\ \delta=0 \quad\text{(R1)}\\
-q^{-1} & \gamma>0,\ \delta=0 \quad\text{(R2)}\\
q & \gamma=0,\ \delta>0 \quad\text{(R3)}\\
1 & \gamma>0,\ \delta>0 \quad\text{(R4)},
\end{cases}
\]
the shifts $z$ running over $\gamma+Nm$, $\delta+Nm$, $\gamma+\delta+N(m-1)$ and $Nm$ as
the rule dictates, and every sum terminating at the first $z$ for which the displayed
wedge fails to be ordered. We call $\varepsilon\,u_{k_2}\wedge u_{k_1}$ the
\emph{transposition term} and the rest the \emph{corrections}.
\end{definition}

\section{The defect in the level-$\ell$ straightening}\label{sec:bug}

Our straightener bubble-sorts, resolving the first inversion by
Definition~\ref{def:R}, and contained the line
\begin{quote}
\texttt{if len(set(idx)) < len(idx): continue \quad\# (R1): $u_k\wedge u_k=0$}
\end{quote}
which discards a wedge as soon as \emph{some} index is repeated, anywhere. (R1) justifies
that only for an \emph{adjacent} repeat. At $\ell=1$ every pair has $\delta=0$, so only
(R1) and (R2) fire, both with $\varepsilon\in\{-1,-q^{-1}\}$: the leading behaviour is
antisymmetric and repeated indices really do die. At $\ell\ge2$, rule (R4) has
$\varepsilon=+1$ --- symmetric --- and they do not. A second, latent defect sat in the same
place: the inversion test \texttt{idx[i] < idx[i+1]} is strict, so a tuple with adjacent
\emph{equal} entries has ``no inversion'' and, once the shortcut is removed, would be
accepted as an ordered wedge. Equality must be handled explicitly, and it is exactly the
case (R1) does cover.

\begin{remark}[The detector that found it, and why the old ones could not]
\label{rem:detector}
The implementation was validated at level $1$ against an independently generated answer
key ($201$ blocks, $818$ coefficients, $0$ mismatches). That check cannot see this bug:
it never fires (R3) or (R4). The check that does is internal and untuned ---
$[B_{-1},B_{-2}]=0$, from Uglov's proposition that $[B_m,B_{m'}]=0$ unless $m+m'=0$.
Both operators raise $k$, so a truncated wedge has no bottom-boundary artefact, and the
test needs neither the multipartition dictionary nor the Chevalley action. On eight
configurations it reported: $7$ nonzero commutators at $\ell\ge2$, $0$ at $\ell=1$. After
the repair: $0$ of $8$, and confluence (straighten resolving the first inversion versus
the last) $8$ of $8$ both before and after --- so confluence alone would not have caught
it either.

This is the fifth consecutive session in which the failure was in the evidence layer and
not the claim layer, and the second in which a green detector was green because it was
blind. The generalisable form: \emph{a validation suite inherited from a special case
certifies only the code paths that special case executes.} Level $1$ runs two of the four
ordering rules.
\end{remark}

\section{$[e_i,B_{-1}]=0$, and the normalisation is trivial}\label{sec:comm}

\begin{theorem}[Gap 3, settled]\label{thm:norm}
Under the indexation of \S\ref{sec:conv}, $u_{\bk}=\ket{\bl,\bs}$ with no normalisation
factor, and $[e_i,B_{-1}]=[f_i,B_{-1}]=0$.
\end{theorem}

\begin{proof}
This is Uglov's, quoted in \S1: the wedge action written in the multipartition indexation
is \emph{identical} to the combinatorial action of Definition~\ref{def:chev}, and the
actions of $U'_q(\widehat{\mathfrak{sl}}_n)$ and $\mathcal H$ commute. No normalisation
is available to be non-trivial, since the two bases are declared equal.
\end{proof}

\begin{remark}[Verification, with the controls]\label{rem:verif-comm}
Recomputed with the repaired straightening over
$(n,\ell,\bs)$ in eleven settings --- $\ell=2$ with $n=2,3,4$ and charges
$(0,0),(0,1),(0,2),(1,0)$; $\ell=3$ with $n=2,3$ --- all multipartitions up to size $4$
(size $2$ at $(3,3)$) and all $i$: $[e_i,B_{-1}]=0$ on $523$ of $523$.
For $f_i$: $515$ of $523$, and all eight exceptions are truncation artefacts --- each is a
single column whose target has at least as many rows as the stored bead depth, and each
vanishes once the depth is raised (all eight are $0$ at depth $\ge9$). This asymmetry is
expected: $f_i$ grows the partition and $e_i$ shrinks it.

Three controls, all of which fire:
\begin{enumerate}
\item \emph{Mirror conventions.} Reversing the tie-break gives $20$ of $96$ commuting;
reversing the side gives $15$ of $96$; reversing both, $13$ of $96$. So the test
discriminates the convention, and Definition~\ref{def:chev} is pinned by it.
\item \emph{Injected normalisation.} Rescaling one basis vector,
$\ket{((1),(1))}\mapsto q\ket{((1),(1))}$, and conjugating $B_{-1}$ by it: the test
flags $5$ cases and exactly the $5$ whose $B_{-1}$-matrix meets that vector, leaving $31$
clean. A detector that cannot see a planted normalisation verifies nothing about the
absence of one.
\item \emph{Charge leakage.} The straightening corrections can move a bead between
runners, which would leave the sector $\mathcal F_{\bs}$; such terms were counted rather
than silently dropped. Their total is $0$, as Uglov's weight argument requires.
\end{enumerate}
\end{remark}

\section{The closed form}\label{sec:closed}

Gap 3 also recorded a measurement: on single-component entries the ratio of $B_{-1}$ to
the naive per-runner ribbon operator ``was always a power $q^{a}$ with $0\le a\le\ell-1$,
but not constant on a runner''. That measurement survives the repair --- and now has a
proof and an identification.

\begin{lemma}[Confinement of a single exchange]\label{lem:conf}
In every term of Definition~\ref{def:R}, both indices lie in the closed interval
$[k_1,k_2]$. Moreover, writing a correction as $u_{k_2-z}\wedge u_{k_1+z}$ with $z>0$, we
have $k_1<k_1+z<k_2-z<k_2$.
\end{lemma}

\begin{proof}
The transposition term has index multiset $\{k_1,k_2\}$. Every correction has $z>0$
(inspect the four rules: the shifts are $\gamma+Nm$ and $Nm\,(m\ge1)$ in (R2);
$\delta+Nm$ and $Nm\,(m\ge1)$ in (R3); $\delta+Nm$, $\gamma+Nm$,
$\gamma+\delta+N(m-1)\,(m\ge1)$ and $Nm\,(m\ge1)$ in (R4); all are positive because
$\gamma,\delta>0$ wherever they appear). The stated cut-off --- summations run only while
the displayed wedge remains ordered --- is precisely $k_2-z>k_1+z$.
\end{proof}

\begin{lemma}[Locality of $B_{-1}$]\label{lem:local}
Fix $j$ and put $p=k_j$, $x=p+N$, $I=[p,x]\cap\Z$. Let $T_j$ be the wedge obtained from
$u_{\bk}$ by replacing $k_j$ with $x$. Then every ordered wedge occurring in the
straightening of $T_j$ has index multiset $\mathbf m$ with
$\mathbf m\setminus I=M\setminus I$ as multisets.
\end{lemma}

\begin{proof}
We show by induction on the straightening steps that at every stage the word of indices
satisfies: (i) the entries $>x$ form a strictly decreasing prefix, and the entries $<p$ a
strictly decreasing suffix, each equal as a sequence to the corresponding part of $M$;
(ii) all remaining entries lie in $(p,x]$.

For $T_j$ this holds: the entries of $u_{\bk}$ greater than $x$ are among
$k_1,\dots,k_{j-1}$ and occupy the prefix; those less than $p$ are $k_{j+1},\dots$ and
occupy the suffix; the middle consists of $x$ together with the $k_r\in(p,x)$, and $p$
itself is gone.

For the step, let $(a,b)$ be the adjacent pair at the first inversion, so $a<b$. If
$b<p$ then $a<b<p$, so both lie in the suffix, which is strictly decreasing --- no
inversion there. If $a>x$ then $b>a>x$, so both lie in the prefix, likewise strictly
decreasing. Hence $a,b\in(p,x]$, and by Lemma~\ref{lem:conf} the two replacement entries
lie in $[a,b]\subseteq(p,x]$. Properties (i) and (ii) are preserved, and no entry outside
$I$ was created or destroyed.
\end{proof}

\begin{proposition}[Leading/correction dichotomy]\label{prop:dich}
Write $M^{(i)}=M\setminus\{k_i\}\cup\{k_i+N\}$. For $i\neq j$, the straightening of $T_j$
contributes nothing to the ordered wedge on $M^{(i)}$. Its contribution to the ordered
wedge on $M^{(j)}$ uses transposition terms only.
\end{proposition}

\begin{proof}
$M^{(i)}$ differs from $M$ exactly at the two values $k_i$ (absent) and $k_i+N$
(present). By Lemma~\ref{lem:local} a term of $B_{-1}$ arising from $T_j$ agrees with $M$
off $I$, so both $k_i$ and $k_i+N$ must lie in $I$. As $I$ has length exactly $N$, this
forces $k_i\ge p$ and $k_i+N\le p+N$, i.e.\ $k_i=p=k_j$ and $i=j$.

For $i=j$: put $\Phi(\mathbf m)=\sum_r m_r^2$. A transposition term leaves the multiset
unchanged, so $\Phi$ is constant across it. A correction replaces $\{k_1,k_2\}$ by
$\{k_1+z,k_2-z\}$ and changes $\Phi$ by $2z\,(k_1-k_2+z)$, which is
$<0$ because $k_2-k_1>2z$ by Lemma~\ref{lem:conf}. Hence $\Phi$ is non-increasing along
any straightening path and strictly decreasing at every correction. Since
$\Phi(M^{(j)})=\Phi(T_j)$, any path from $T_j$ to $M^{(j)}$ used no correction.
\end{proof}

\begin{theorem}[Single-component entries of $B_{-1}$]\label{thm:closed}
Let $\nu\neq\bl$ with $\bra{\nu}B_{-1}\ket{\bl}\neq0$ and $\nu,\bl$ differing in a single
component $d$. Then $M_d(\nu)=M_d(\bl)\setminus\{\kappa\}\cup\{\kappa+e\}$ for a unique
$\kappa$ with $\kappa\in M_d(\bl)$, $\kappa+e\notin M_d(\bl)$, and
\[
\boxed{\ \bra{\nu}B_{-1}\ket{\bl}\;=\;(-q^{-1})^{h}\,q^{a},\qquad
h=\#\big(M_d\cap(\kappa,\kappa+e)\big),\ }
\]
\[
a\;=\;\#\{d'>d:\ \kappa\in M_{d'}\}\;+\;\#\{d'<d:\ \kappa+e\in M_{d'}\}.
\]
\end{theorem}

\begin{proof}
By Proposition~\ref{prop:dich} the entry receives contributions only from $T_j$ with
$k_j=p$ the $k$-coordinate of $(d,\kappa)$, and only through transposition terms. So it
is the product of the leading coefficients $\varepsilon$ collected while carrying
$u_{p+N}$ leftwards past every bead $k_r$ with $p<k_r<p+N$; there are no other steps,
since after passing them the word is ordered.

Fix such a $k_r$ and apply Definition~\ref{def:R} to $(k_1,k_2)=(k_r,p+N)$. The label of
$p+N$ is the label $(c,d)$ of $p$.
\begin{itemize}
\item If $d_r=d$ then $\delta=0$. We cannot have $c_r=c$ as well: that would make
$k_r\equiv p \pmod N$, impossible for $p<k_r<p+N$. So $\gamma>0$, rule (R2),
$\varepsilon=-q^{-1}$.
\item If $d_r\neq d$ then $\delta\neq0$, since $0<|n(d_r-d)|<N$. If $c_r=c$ we are in
(R3) and $\varepsilon=q$; otherwise (R4) and $\varepsilon=1$.
\end{itemize}
Thus the product is $(-q^{-1})^{h'}q^{a'}$ with $h'$ the number of beads in the window on
runner $d$ and $a'$ the number on other runners with the same $c$.

Both counts are the stated ones. In the $\kappa$-coordinate the window
$(p,p+N)$ meets runner $d$ in $(\kappa,\kappa+e)$, so $h'=h$. For $a'$: the beads in the
window with $c$-label $c$ are those with label $(c,d')$, $d'\neq d$, and by the second
remark of \S\ref{sec:conv} each such label occurs exactly once in the window. Solving
\eqref{eq:kcd} for that occurrence, $k'=c+n(d'-1)-N\mu'$ lies in $(p,p+N)$ iff
$0<n(d'-d)-N(\mu'-\mu)<N$, which gives $\mu'=\mu$ when $d'>d$ and $\mu'=\mu-1$ when
$d'<d$; in the $\kappa$-coordinate, $\kappa'=\kappa$ and $\kappa'=\kappa+n$ respectively.
So $a'$ counts $d'>d$ occupied at $\kappa$ and $d'<d$ occupied at $\kappa+e$, which is
$a$.
\end{proof}

\begin{corollary}\label{cor:range}
$0\le a\le\ell-1$, and $0\le h\le e-1$.
\end{corollary}

\begin{proof}
Each label other than $(c,d)$ occurs exactly once in the window: the labels contributing
to $a$ are the $\ell-1$ labels $(c,d')$ with $d'\neq d$, and those contributing to $h$
are the $n-1$ labels $(c',d)$ with $c'\neq c$.
\end{proof}

\begin{corollary}[$a$ is a cross-runner quantity]\label{cor:cross}
$a$ depends only on the configuration on the runners $d'\neq d$, at the two positions
$\kappa$ and $\kappa+e$. In particular it is not determined by runner $d$, so it is not
constant along a runner and cannot be absorbed by any normalisation that is a function of
the component sizes.
\end{corollary}

\begin{remark}[Verification]\label{rem:verif-closed}
Over twelve settings $(n,\ell,\bs)$ --- $\ell=2$ with $n=2,3,4$ and four charges,
$\ell=3$ with $n=2,3$ --- and all multipartitions up to size $3$: the boxed formula
reproduced \emph{every} single-component entry of $B_{-1}$, $913$ of $913$, both as a
ratio-exponent and as the full coefficient. The distribution of $a$ is not concentrated:
at $\ell=2$ it is $\{0,1\}$ in the proportion $84{:}84$ and at $\ell=3$ it is
$\{0,1,2\}$ in the proportion $56{:}56{:}56$, and within a fixed runner $d$ all values
occur (e.g.\ $\ell=3,d=1$: $a=2$ sixty-eight times, $a=1$ nineteen, $a=0$ eight).
Separately, the dichotomy of Proposition~\ref{prop:dich} was measured directly by tagging
each straightening path with whether it used a correction: single-component entries
received $697$ leading-only contributions and $0$ correction contributions;
two-component entries received $0$ and $118$. There is no cancellation --- the two
classes of path have disjoint targets, exactly as the locality lemma predicts.
\end{remark}

\section{Two corollaries promoted to \texttt{proved}}\label{sec:promote}

The 31 August paper recorded these as computational observations inside
Proposition~\ref{prop:heis-old}. They follow from Proposition~\ref{prop:dich}.

\begin{proposition}[31 Aug., Prop.\ 4.1]\label{prop:heis-old}
For $\ell\ge2$, $B_{-1}$ has nonzero entries between multipartitions differing in two
components; hence no single-bead-move operator specialises to it.
\end{proposition}

\noindent The existence statement is computational and re-verified after the repair
($55$ two-component entries in the narrow sweep, $118$ in the wide one), and the support
argument in the 31 August paper is unaffected. What is new:

\begin{corollary}\label{cor:div}
Every entry $\bra{\nu}B_{-1}\ket{\bl}$ with $\nu$ and $\bl$ differing in two or more
components is divisible by $(q-q^{-1})$ in $\Z[q,q^{-1}]$.
\end{corollary}

\begin{proof}
By Proposition~\ref{prop:dich} the transposition-only paths from $T_j$ land exactly on
$M^{(j)}$, which differs from $M$ in one component. So such an entry receives only paths
using at least one correction, and its value is a sum of products each containing at
least one correction coefficient. Every correction coefficient is $(q-q^{-1})$ times a
Laurent polynomial: in (R2), $(q^{-2}-1)q^{-2m}=-q^{-1-2m}(q-q^{-1})$ and
$(q^{-2}-1)q^{-2m+1}=-q^{-2m}(q-q^{-1})$; in (R3), $(q^{2}-1)q^{\pm}$ likewise; in (R4)
the factor $(q-q^{-1})$ is explicit, and the cofactors
$\frac{q^{2m+1}+q^{-2m-1}}{q+q^{-1}}=\sum_{i=0}^{2m}(-1)^{i}q^{2m-2i}$ and
$\frac{q^{2m}-q^{-2m}}{q+q^{-1}}=\sum_{i=0}^{2m-1}(-1)^{i}q^{2m-1-2i}$
are Laurent polynomials. The transposition coefficients are units.
\end{proof}

\begin{corollary}\label{cor:q1}
At $q=1$, $B_{-1}$ is exactly the naive per-runner ribbon operator: for every $\bl$,
\[
B_{-1}\big|_{q=1}\ket{\bl}\;=\;\sum_{d}\ \sum_{\substack{\kappa\in M_d,\ \kappa+e\notin M_d}}
(-1)^{h}\,\ket{\bl^{(d,\kappa)}} .
\]
\end{corollary}

\begin{proof}
By Corollary~\ref{cor:div} all multi-component entries vanish at $q=1$. By
Theorem~\ref{thm:closed} the single-component entries are $(-1)^{h}q^{a-h}$, which is
$(-1)^{h}$ at $q=1$.
\end{proof}

\noindent (Both were previously supported by counts --- $118$ divisibility checks and
$487$ entrywise $q=1$ agreements, all of which were recomputed after the repair and all
of which still pass.)

\section{(H4): refuted, and what it got right}\label{sec:h4}

\begin{hypothesis}[(H4), from the brief]
The unidentified $q$-power is the tie term: the correct identification is
$u_{\bk}\leftrightarrow q^{\tau(\bl,\bs)}\ket{\bl,\bs}$ up to a global constant, with
$\tau$ the cross-runner tie term of the 31 August telescope theorem.
\end{hypothesis}

\noindent \textbf{Refuted}, together with its named rival $a=\#\{\text{occupied
runners}\}-1$ and with any other candidate, because there is no normalisation power to
identify: Theorem~\ref{thm:norm}. The premise was an artefact.

The brief anticipated that the data might fail to separate (H4) from its rival and called
that ``a finding, and the honest one''. The actual outcome is sharper and it is worth
naming the difference: \emph{the rival test was never reached, because the quantity both
hypotheses were about did not exist.} The brief's four supporting coincidences --- matching
ranges, both non-constant on a runner, both vanishing at $\ell=1$, and the prospect of
unifying two gaps --- were all consistent with the truth and all pointed at the wrong
object. The one the brief rated strongest, ``both are not constant on a runner'', is
exactly the one that survives: Corollary~\ref{cor:cross} proves non-constancy, for the
cross-runner reason (H4) guessed.

So (H4) had the right species and the wrong home. The $q$-power is a cross-runner count
at tied positions, bounded by $\ell-1$, vanishing at $\ell=1$ --- every property (H4)
predicted. It is simply a feature of the operator $B_{-1}$ rather than of the basis. A
hypothesis can be refuted and still have been the right thing to write down; what made
this cheap was stating it precisely enough that ``is there such a power at all?'' became
the first question rather than the last.

\begin{remark}[Where the four days of trouble actually lived]
Three sessions treated the wedge/Fock identification as the suspect object. It was never
the suspect; the suspect was six lines of straightening code whose only validation came
from a level at which two of its four branches are unreachable. The record of the
programme said ``normalisation'' because the first session to see the symptom wrote that
down, and the next two read it. This is the sixth instance of
\emph{recorded facts calcify}, and the specific corrective is the one that worked today:
before hunting a correction factor, check whether the source says the factor is $1$.
\end{remark}

\section{Consequence for the object of Q59}\label{sec:consequence}

The 31 August paper closed by naming its main open problem: find the level-$\ell$ operator
whose specialisations are $P_e$ and $B_{-1}$, ``presumably by adding the
$(q-q^{-1})$-divisible cross-runner terms with a second grading''. Theorem~\ref{thm:closed}
supplies the single-bead half of that operator exactly. Where the 31 August graded ribbon
operator carries $t^{h}$ on the move $(d,\kappa)\mapsto(d,\kappa+e)$, the operator that
interpolates must carry $q^{a}t^{h}$ with $a$ as in Theorem~\ref{thm:closed}: at $t=-q^{-1}$
this reproduces $B_{-1}$ on every single-component entry, by the theorem. What remains
unpinned is the cross-runner part, which by Corollary~\ref{cor:div} is $(q-q^{-1})$-divisible
and therefore invisible at $q=1$.

This is a statement about what to try next, not a result. In particular it is \emph{not}
claimed that reweighting by $q^{a}$ repairs Theorem 3.1 of the 31 August paper (the
$(1+q^{j}t)$ dichotomy): that theorem is about the unweighted operator, its evidence is
unaffected by anything here --- it uses no wedge realisation and no straightening --- and
whether the reweighted operator satisfies a rigidity statement at all is untested.

\section{Gaps}\label{sec:gaps}

\begin{enumerate}
\item \textbf{The cross-runner entries have no closed form.} Corollary~\ref{cor:div}
bounds them ($(q-q^{-1})$-divisible, arising only from correction paths) but does not
compute them. This is now the sharp remaining question about $B_{-1}$ at level $\ell$, and
Lemma~\ref{lem:local} localises it: such an entry is determined by the beads in a single
window of length $N$.

\item \textbf{Theorem~\ref{thm:norm} is cited, not reproved.} Uglov's assertion that the
two actions are identical is stated without proof at that point in his paper (it follows
from his construction). We verified it computationally rather than deriving it. That is
the right weight to give it --- but the node's trust should read \texttt{cited} plus
\texttt{computed}, not \texttt{proved}.

\item \textbf{Sizes are small.} All sweeps are $|\bl|\le4$, most $\le3$, with bead depths
$5$--$9$. The $f_i$ truncation artefacts of Remark~\ref{rem:verif-comm} are a reminder
that the depth is a real parameter and that the honest reading of a failure is
``check the depth first''.

\item \textbf{Conjecture \texttt{conj:j} was not touched.} It was step 2 of the brief and
the session did not reach it. It is undamaged by anything here: it concerns
$[e_i,R^{(\ell)}(t)]$, which uses neither the wedge nor the straightening.
\end{enumerate}

\section{Corrections to the 31 August paper}\label{sec:corrections}

\begin{enumerate}
\item \S Gaps item 3 (``the wedge/Fock normalisation is unresolved \dots\ the
identification carries an unidentified $q$-power'') is \textbf{withdrawn}. The
normalisation is trivial; the $13/126$ was the straightening defect of \S\ref{sec:bug}.
\item Remark 1.4 (``the tie-break direction and the choice of side \dots\ were fixed by
requiring that the $\ell=1$ specialisation reproduce the proved Q59 theorem'') is
\textbf{corrected}: the $\ell=1$ reduction pins the side and cannot see the tie-break.
The tie-break is pinned by Uglov's Theorem 2.1 and, independently, by
$[e_i,B_{-1}]=0$ (Remark~\ref{rem:verif-comm}).
\item Proposition 4.1 and its verification remark \textbf{stand}, recomputed after the
repair; two of its computational claims are upgraded to proofs
(Corollaries~\ref{cor:div}, \ref{cor:q1}).
\item Theorems 2.3 and 3.1, Corollary 2.4 and Conjecture 3.3 are \textbf{unaffected}:
none of them uses the wedge realisation or the straightening.
\end{enumerate}

\end{document}
